\section{Translated characters of higher order}\label{sec:characters}

In this section the harmonic mass of a set of primes \(\mathcal P\) is
\(\sum_{p\in\mathcal P}p^{-1}\). All multiplicative characters are
extended by zero at zero. The parameters introduced in this section
are local to its proof.

We now prove translated decorrelation for every character of order
greater than two. Repeated Cauchy--Schwarz steps create copies of
the prime variables, and pairs of smaller and larger anchor primes
distinguish their positions. Changing a position then exposes a
one-sided interaction involving the square of a selected character.
That square is nonprincipal because the character has order greater
than two. The final corollary combines this argument with the
quadratic conclusion of Section~\ref{sec:quadratic}.

\begin{proposition}\label{prop:characters}
For every fixed \(0<\alpha<\beta\), one has
\begin{equation}\label{eq:src16}
 \sum_{\alpha L\le \log\log p\le\beta L}\frac1p
 \max_{\substack{t\in\F_p,\ \lambda\ {\rm multiplicative}\\
                  \ord(\lambda)>2}}
 \left|\E_{x\in S_p}\lambda(x-t)\right|=o(L)
 \qquad(L\longrightarrow\infty).
\end{equation}
\end{proposition}

\begin{proof}
Suppose otherwise. After passing to a sequence \(L\to\infty\),
there are constants \(c,\delta>0\), sets \(\mathcal E\) of primes in
the indicated band with harmonic mass at least \(cL\), and choices
\[
 f_p(x)=z_p\chi^{(p)}(x-t_p),\qquad
 |z_p|=1,\qquad \ord(\chi^{(p)})>2,
 \qquad \Re\E_{x\in S_p}f_p(x)\ge\delta .
\]
All constants below may depend on \(\alpha,\beta,c,\delta\).
Whenever labels are sampled from a specified prime set, their prior
is the harmonic measure on that set, normalized to a probability.
Restrictions on a tuple, such as distinctness or a product bin,
will be imposed by indicators, without renormalizing its prior.

\pdfbookmark[2]{Selection of scales and positive statistics}{characters.selection}
\subsection*{Selection of scales and positive statistics}

Set \(\epsilon=10^{-4}\). We shall choose a sufficiently large
constant \(B_D\), followed by a sufficiently large lower bound \(K_0\)
on an integer depth \(k\). Leading constants multiplying \(m\) in
the estimates below, except for the explicitly displayed
\(\log z\) terms, will be independent of \(B_D,k\). Symbols such as
\(O_k(1)\) and \(o_k(m)\) allow dependence on these fixed choices;
their limiting variable is \(L\).

Use the coordinate \(u=\log\log p\). Mertens' estimate, uniformly
on the intervals in question, gives harmonic mass at most the
interval length plus \(o(1)\). We can therefore find three
subintervals in increasing order, each containing \(\gg L\) mass
of \(\mathcal E\), separated from each other by gaps \(\gg L\).
Indeed, partition the original band into \(q\) equal intervals,
where \(q\) is a sufficiently large fixed integer. Intervals of
mass less than \(cL/(2q)\) carry together less than \(cL/2\).
Since each other interval has mass at most
\((\beta-\alpha)L/q+o(1)\), there are at least
\(cq/(3(\beta-\alpha))\) other intervals for large \(L\).
Take \(q\) large enough that this number exceeds 12, and select
three of these intervals with at least one whole partition
interval between successive choices. Their lengths and gaps
are fixed positive multiples of \(L\).
We use all the \(\mathcal E\)-primes in the middle interval as
\emph{bulk labels}. In the lower interval, a unit-length shell
has \(\mathcal E\)-mass bounded below by a positive constant;
fix one such shell for the small anchors.

We require more structure in the upper interval. There is a fixed
\(c_1>0\) and a block \([U,U+5k]\), with \(k\ge K_0\) bounded above
independently of \(L\), such that every subinterval of length
\(\epsilon k\) in the block has \(\mathcal E\)-mass at least
\(c_1\epsilon k\). Here is the density-increment argument, including
the dependence of the constants. Choose an integer
\(M>100/\epsilon\). A block of fixed length \(5K_0M^h\) inherits
some fixed positive density \(d\) from the upper interval, by
averaging over a partition into such blocks. If a current block of
length \(5k\) has density at least \(d\) but contains a failing
interval of length \(\epsilon k\), the child blocks of length
\(5k/M\) wholly contained in that interval have total length at
least \(\epsilon k/2\). Choose \(c_1<d/4\). These children have
average density at most \(2c_1<d/2\), so one of the other children
has density larger than the current density by a fixed positive
amount depending only on \(d,\epsilon\). This increase cannot occur
more than \(h\) times, for sufficiently large fixed \(h\), because
all densities are at most \(1+o(1)\). Neither \(h\) nor \(c_1\)
depends on \(K_0\). Thus the terminal \(k\) belongs to the finite set
\(\{K_0,K_0M,\ldots,K_0M^h\}\). Pass to an unbounded subsequence
on which \(k\) is fixed.

Increase \(K_0\) so that discarding at most two boundary unit
shells from an interval of length \(\epsilon k\) costs less than
half its guaranteed mass. Every such interval then contains a
whole shell from the grid \(U+\Z\) with \(\mathcal E\)-mass at
least a fixed \(c_0>0\). Call these shells \emph{rich}. Fix one
rich shell in \([U,U+2\epsilon k]\); it supplies a top label and
the big anchors. Let \(\tau\) be the exponential of its lower
endpoint in the \(u\)-coordinate, and set
\[
 \log X=10^3\,4^k\tau,\qquad
 z=\exp(3\epsilon k),\qquad m=\lfloor zL\rfloor .
\]

We next obtain good endpoints simultaneously for every shell that
may be needed. The sets to be tested are the bulk interval, the
small-anchor shell, the top shell, and all rich grid shells in the
late block lying wholly below \(\log p=(\log X)/4\). Their number
is bounded for fixed \(k\), and each has harmonic mass bounded
below by a positive constant. Let
\(\mathcal A_X=A\cap[X^{9/10},X]\).
For any subset \(\mathcal A'\subset\mathcal A_X\) with
\[
 |\mathcal A'|\ge\frac{\sqrt X}{(\log X)^{10}},
\]
apply Lemma~\ref{lem:collision} to its uniform measure and the
uniform measure on the full \(D\)-tail. Every tested prime is
below the cutoff in that lemma. Cauchy--Schwarz on the residue
classes, followed by \(\log p\ge e^{\alpha L}\), gives
\[
 \sum_{\text{tested }p}\frac1p
 \left|\E_{a\in\mathcal A'}f_p(a)
       -\E_{x\in S_p}f_p(x)\right|^2
 \ll (\log\log X)e^{-\alpha L}=o(1).
\]
If at least \(\sqrt X/(\log X)^{10}\) endpoints failed
\(\Re\E_p f_p(a)\ge\delta/2\) for one of the tested prime sets,
their uniform measure would contradict this estimate and
Cauchy--Schwarz over that set of primes. Lemma~\ref{lem:size}
shows that \(|\mathcal A_X|\gg\sqrt X/(\log X)^3\).
Consequently, outside a negligible fraction of \(\mathcal A_X\),
all these tests hold simultaneously. Call these endpoints typical.

For a rich unit shell with lower exponential scale \(b\), partition
its primes into the cells
\[
 \mathcal E_h=\{p\in\mathcal E:h\le\log p<h+1\},
 \qquad h\in\Z\cap[b,eb].
\]
Boundary cells have negligible mass, and the prime-counting upper
bound gives mass \(O(1/b)\) for every cell. At a typical endpoint,
the shell average has real part at least \(\delta/2\). Since
\(|f_p|\le1\), a fixed positive amount of shell mass lies in cells
whose average has real part at least \(\delta/4\). Discarding cells
with mass smaller than a sufficiently small constant times \(1/b\)
loses less than half this amount. Thus there are \(\gg b\) good
cells satisfying
\[
 \sum_{p\in\mathcal E_h}\frac1p\gg\frac1b,
 \qquad
 \Re\E_{p\in\mathcal E_h}f_p(a)\ge\delta/4.
\]
The constants are uniform over every tested shell.

A word consists of \(m\) independent bulk labels and one independent
top label. Define
\[
 G_J(n)=\E_{\rm word}
  \ind_{\{\lfloor\sum\log p\rfloor=J\}}\prod f_p(n).
\]
There are at most \(\exp(CL)\) possible bins. The gap between the
bulk and top scales implies
\[
 (1-o(1))\tau\le J\le4\tau
\]
whenever the bin is nonempty. At a typical endpoint,
independence gives
\(\big|\sum_JG_J(a)\big|\ge(\delta/2)^{m+1}\).
Pigeonholing first a bin for each endpoint and then the endpoints
among bins yields one fixed \(J\) such that
\[
 |G_J(a)|\ge e^{-Cm}
\]
at at least \(\sqrt X e^{-Cm}\) typical endpoints. The constants
here are independent of large fixed \(k\): the bin count and the
tail-size logarithmic loss have logarithms \(O(L)+O_k(1)\).

Introduce the gaps and frequency bounds
\begin{align}
 \Delta_0&=(B_D+20\log z)m,\qquad
 \Delta_j=2^{j-1}\Delta_0+4^j\sqrt m
       &&(1\le j\le k),\notag\\
 E_j&=2^j\Delta_0+2\cdot4^j\sqrt m,\qquad
 V_j=e^{E_j}
       &&(0\le j\le k).
 \label{eq:src17}
\end{align}
For each retained endpoint choose a good small-anchor cell and a
good big-anchor cell for every \(j\), and let \(a_j\) be the sum
of their two indices. We shall choose a prime group \(K_j\) of
logarithmic target \(T_j\) to be removed at step \(j\), and a
filler group of target \(T_F\). At step \(j\), the part to be
copied will consist of \(2^{j-1}\) copies of the selected word,
one group of each future pivot type, and the two fresh anchors.
We require its logarithmic size to exceed \(T_j\) by \(\Delta_j\).
The filler then makes the full product in the initial squared
expression have logarithmic size \(\log X+\Delta_0\), as needed
for Poisson summation. These requirements give the backwards definitions
\[
 T_j=2^{j-1}J+\sum_{\ell>j}T_\ell+a_j-\Delta_j,
 \qquad
 T_F=\frac{\log X+\Delta_0}{2}
       -J-\sum_{j=1}^k(T_j+a_j).
\]
They are positive and, for all sufficiently large \(L\),
\[
 c_*\,2^k\tau\le T_j,T_F\le600\,4^k\tau
\]
with an absolute \(c_*>0\). To check this, put
\(c_j=2^{j-1}J+a_j-\Delta_j\). The backwards recurrence gives
\[
 \sum_jT_j=\sum_j2^{j-1}c_j.
\]
Here \(a_j\le4\tau\) eventually and \(\Delta_j=o_k(\tau)\);
the contribution of \(J\) is
\((4^k-1)J/3\), while the anchor contribution is \(O(2^k\tau)\).
These estimates prove the bounds on the \(T_j\), and the large
coefficient \(10^3\) proves the bounds on \(T_F\).

For any such target \(T\), the interval
\[
 [\log T-2\epsilon k,\log T-\epsilon k]
\]
lies inside the late block when \(K_0\) is large. It contains a
rich unit shell. The primes of that shell obey
\(\log p\le T e^{-\epsilon k}<(\log X)/4\), again by increasing
\(K_0\); hence this shell was among the simultaneous endpoint
tests above.

We record why the good cells in that shell can realize the target
to bounded accuracy. Let their index set be \(I\subset[b,eb]\),
with \(|I|\ge c'b\), and write
\[
 h_-=\min I,\quad h_+=\max I,\quad s=h_+-h_-,
 \quad g=\gcd(I-h_-).
\]
Then \(s\gg b\) and \(g\ll1\). Let \(\mathcal U\) be the set
\((I-h_-)/g\), reduced modulo \(s/g\). It contains zero, generates
the group, and has density bounded below. A bounded number \(r_0\)
of its sums covers that group. For completeness, Kneser's sumset
theorem \cite{Kneser1953}, in the form of
\cite[Theorem~1]{DeVos}, gives, for the stabilizer \(H\) of an
\(r_0\)-fold sumset,
\[
 |r_0\mathcal U|
       \ge r_0|\mathcal U+H|-(r_0-1)|H|.
\]
If \(H\) is proper, the generating set \(\mathcal U\) meets at least
two \(H\)-cosets, so the right side is at least \(r_0|\mathcal U|/2\).
For sufficiently large fixed \(r_0\) this is impossible in the
ambient group. Thus the sumset is the whole group. Appending zero
and endpoint summands now shows that \(n\) sums of \(I-h_-\)
cover every multiple of \(g\) in
\([r_0s,(n-r_0)s]\).
Choose \(n\) to be the integer nearest
\(T/((h_-+h_+)/2)\). Since the shell lies in the preceding target
interval,
\[
 c''e^{\epsilon k}\le n\le C e^{2\epsilon k}.
\]
For large \(K_0\), the target lies in the interior interval just
described. Rounding to the appropriate residue class modulo \(g\)
therefore gives an ordered list of \(n\) good cell indices with
sum \(T+O(1)\).

Use such lists for every pivot target \(T_j\) and the filler target
\(T_F\). The total number \(n_c\) of cell labels, including the
anchors, satisfies
\[
 n_c\le n_{\max}\ll k e^{2\epsilon k}.
\]
Every index has at most \(\exp(CL)\) possibilities. Fixing all
lengths, ordered lists, and anchor choices by pigeonhole costs at
most
\[
 \exp(Cn_{\max}L+O_k(1)).
\]
Since \(n_{\max}/z\ll k e^{-\epsilon k}\to0\), this is absorbed
by \(\exp(Cm)\), with a leading constant independent of \(k,B_D\).
The endpoint dependence of the targets causes no loss beyond this
pigeonhole: every eligible rich shell had already been tested at
every typical endpoint.

Fix a real nonnegative Schwartz function \(\psi\), bounded below
by a positive constant on \([0,1]\), with smooth compactly
supported Fourier transform. Such a function is obtained by
squaring the real inverse Fourier transform of a smooth even
nonnegative bump sufficiently concentrated near zero.
For a selected cell write
\(R_h(n)=\E_{p\in\mathcal E_h}f_p(n)\). The selected endpoints
give the positive statistic
\begin{equation}\label{eq:src18}
 I_*=\sum_{n\in\Z}\psi(n/X)|G_J(n)|^2
             \prod_h|R_h(n)|^2
       \ge \sqrt X\,e^{-Cm},
\end{equation}
where cell labels occur with their selected multiplicities.
Its expansion uses an independent positive and conjugate copy
of every role. For the formal product \(M\) of all prime labels,
with multiplicity, the bins and target relations give
\[
 \log M=\log X+\Delta_0+O_k(1).
\]

We may discard tuples with repeated primes. A repeated tuple has
actual period at most \(M/p_{\min}\), and
\(p_{\min}\ge\exp(e^{\alpha L})\). Thus its period is much less
than \(X\); bandlimited Poisson leaves only its complete residue
mean. That mean vanishes if a prime occurs just once, because its
character is nonprincipal. For an all-multiple tuple, the absolute
contribution before its prior is \(O(X)\).
Its point weight is at most
\[
 M^{-1}L^{-2m}\exp(Cm+O_k(1)).
\]
Here the bulk priors supply \(L^{-2m}\); the cell normalizations
cost at most \(\exp(Cn_cL)\), and the top normalizations cost a
bounded factor per top label.
Extract
\[
 M^{-1/2}\le X^{-1/2}e^{-\Delta_0/2+O_k(1)}.
\]
There are \(b_*=2m+O_k(1)\) positions. If there are \(d\le b_*/2\)
distinct primes, summing the remaining \(M^{-1/2}\) costs at most
\[
 \sum_{d\le b_*/2}\frac{d^{b_*}(CL)^d}{d!}
 \le b_*^{b_*}e^{CL};
\]
each distinct prime contributes at least a factor \(1/p\).
The total repeated-prime error is therefore at most
\[
 \sqrt X\exp\!\left(
   -\frac{\Delta_0}{2}+(C+2\log z)m+o_k(m)\right).
\]
Choose \(B_D\) sufficiently large. After Poisson summation on the
distinct contribution and division by \(\sqrt X\), we obtain an
amplitude \(\eta_0\) with
\begin{equation}\label{eq:src19}
 |\eta_0|\ge e^{-B_0m},
\end{equation}
where \(B_0\) is bounded independently of \(B_D,k\).

\pdfbookmark[2]{Templates and the transfer identity}{characters.transfer}
\subsection*{Templates and the transfer identity}

The level-zero list consists of both copies of all selected roles.
Mark the positive word and each positive pivot group \(K_j\)
active. Reserve the two positive anchors for each step \(j\).
Negative copies and unused roles, including the filler, stay
outside. Initially regard a pivot group as one atom whose
constituents are its ordered prime slots. Retain its internal
distinctness indicator in its tuple prior, without normalization.
All other atoms are single prime slots. Pairwise coprimality of
atoms imposes the remaining distinctness conditions. Every copied
atom receives new independent priors with the same internal
restrictions; all priors are probability or subprobability measures.

At step \(j\), partition the current list \(\mathcal I_{j-1}\)
into \(P\sqcup H\sqcup Y\). Here \(P\) is the unique active atom
of type \(K_j\); \(H\) consists of all active words, the unique
active atom of each future type \(K_\ell\), \(\ell>j\), and the
two fresh anchors for step \(j\); \(Y\) is the remainder.
The letters \(P,H,Y\) also denote the corresponding integer
products. Replace this list by
\[
 \mathcal I_j=H^+\sqcup H^-\sqcup Y.
\]
Both copies of every active word remain active. For each future
pivot retain only its positive copy as active; its other copy
becomes outside. Both copies of the fresh anchors become outside.
Thus immediately before step \(j\) there are \(r=2^{j-1}\)
active words, and on their bin supports
\begin{equation}\label{eq:src20}
 \log P=T_j+O_k(1),\qquad
 \log H=T_j+\Delta_j+O_k(1).
\end{equation}

For a prime slot \(i\), let \(p_i\) be its assigned prime and
let \(\lambda_i\) be \(\chi^{(p_i)}\) or its inverse according
to the original positive or conjugate role. This choice is copied
unchanged, including into a negative transfer copy.
For a pairwise distinct list \(\mathcal I\), with prime product
\(M_{\mathcal I}\), define
\[
 A_{\mathcal I}(v)=
   \prod_{i\in\mathcal I}
    e_{p_i}\!\left(t_{p_i}v
             \overline{M_{\mathcal I}/p_i}\right).
\]
The bar denotes inversion in \(\F_{p_i}\). Every term will have
phase of the form
\begin{equation}\label{eq:src21}
 \Theta=A_{\mathcal I}(v)\prod_i\nu_i
       \prod_i\prod_{h\ne i}\lambda_i(p_h)^{b_{ih}}.
\end{equation}
When all frequency parameters are fixed, \(\nu_i\) is a bounded
unary function of the prime in its role; the exponents \(b_{ih}\)
are determined by the template. This expression is used only on
coprime support.

At level zero all \(b_{ih}=1\), and
\(\nu_i=\kappa_i\lambda_i(v)^{-1}\), where \(|\kappa_i|=1\).
Indeed, the local positive Gauss transform at
\(v\overline{M_{\mathcal I}/p_i}\) contributes its normalized
Gauss sum and role multiplier, the translation factor above,
and
\(\lambda_i(v)^{-1}\lambda_i(M_{\mathcal I}/p_i)\).
This is exactly the claimed graph phase. The zero frequency
vanishes. The remaining weight \(W_0\) is the product of the word
bin indicators and
\[
 (X/M)^{1/2}\widehat\psi(vX/M),\qquad 0<|v|\le V_0,
\]
together with indicators of pairwise coprime atoms,
\((v,M)=1\), and a fixed covering interval for
\(\log M=\log X+\Delta_0+O_k(1)\).
The compact Fourier support fits inside \(V_0\) for large \(L\).
Thus
\(\eta_0=\E_{\mathcal I_0}\sum_v W_0\Theta_0\).
The nonphase weight sees each active pivot only through its
total product. Its internal tuple prior remains outside that weight.

At level \(l\), the amplitude is an expectation over
\(\mathcal I_l\) and a sum of terms \(W_l\Theta_l\) over binary
frequency histories, with root \(0<|v|\le V_l\).
The erased pivot at every internal node is specified by the
substitution below. The following properties are maintained:
\begin{enumerate}[label=(\roman*)]
\item
The weight and its support depend on current active pivots only
through their totals; all active word bins hold. Current atoms
are pairwise coprime and are units modulo the absolute root
frequency.
\item
For a slot in an active word or active pivot, called regular,
\[
 b_{ih}=\varepsilon_i\quad(h\ne i),\qquad
 \nu_i=\kappa_i\lambda_i(v)^{-\varepsilon_i},
\]
where \(\varepsilon_i\in\{1,-1\}\) is the product of its transfer
copy signs. The multiplier \(\kappa_i\) depends only on the
role, path, and its prime, not on any frequency.
\item
Each incoming row is constant on the constituent targets of any
active pivot atom.
\end{enumerate}

Consider step \(j=l+1\). Insert into \(W_l\) the range indicators
\eqref{eq:src20} with fixed covering constants, and require that
the pivot be a unit modulo every frequency in its history.
These do not change the actual expectation: all such frequencies
are smaller than the smallest actual prime. Group terms by
\(u=v/H\pmod P\). For a pivot constituent \(p_i\), its outgoing
row and unary give
\[
 \kappa_i
 \lambda_i\!\left((P/p_i)YH/v\right)^{\varepsilon_i}.
\]
Together with its additive factor this depends only on
\(u,Y\), and the pivot tuple. It is bounded by 1 in modulus.
Remove all these pivot factors, and call the remaining
sum and average \(B_u(P,Y)\). The common-column property and
the weight property show that \(B_u\) depends on the pivot
tuple only through its total \(P\). Cauchy--Schwarz first on
the outer priors and then on the residues gives
\begin{equation}\label{eq:src22}
 |\eta_l|^2
 \le \E_{Y,K_j}P\sum_{u\bmod P}|B_u(P,Y)|^2
 \le e^{d_j}\E_Y\sum_P\sum_{u\bmod P}|B_u(P,Y)|^2.
\end{equation}
The final sum is over all positive integers in the first range
of \eqref{eq:src20}. If \(n_j\) is the number of constituents
of \(K_j\), its total has point mass at most
\(e^{d_j}/P\), where
\[
 d_j=Cn_jL+O_k(1).
\]
This follows from the cell mass bounds, with the bounded-for-\(k\)
ordering multiplicity supplied by unique factorization.
On extending \(P\), retain its total-product coprimality and
weight conditions but no internal tuple-prior condition.
No character on the extended integer \(P\) is needed: its
outgoing character rows have already been removed.

Expand the square with two copies \(H_L,H_R\) of \(H\), and
previous root frequencies \(v,w\). Set aside the diagonal
\(vH_R=wH_L\). On the complement the common-residue condition
is equivalent to
\begin{equation}\label{eq:src23}
 P=\frac{vH_R-wH_L}{s},\qquad 0<|s|\le V_j,
\end{equation}
with the indicator that \(P\) is a positive integer in range.
The bound on \(s\) follows from
\(E_{j-1}+\Delta_j+O_k(1)<E_j\).
A prime common to \(H_L,H_R\) would divide \(s\), since it is
coprime to \(P\), which is impossible by the frequency bound.
The inherited supports separate either branch from \(Y\).
All output primes exceed \(V_j\), so they are units modulo
the nonzero \(s\). Hence the new output atoms are pairwise
coprime and units modulo \(s\).
Let \(W_j\) be the product of the left weight and the conjugate
right weight, with these and all inherited indicators and with
the substitution \eqref{eq:src23}.

The additive phases become \(A_{\mathcal I_j}(s)\).
At a left-slot prime use \(v/P=s/H_R\), at a right-slot prime
use \(-w/P=s/H_L\), and at a shared outside prime use
\[
 \frac{v}{PH_L}-\frac{w}{PH_R}
       =\frac{s}{H_LH_R}.
\]
These identities are taken modulo the relevant actual prime,
where all denominators are units.
Write \(b_{iP}\) for the common incoming exponent into the
pivot. For \(t,u'\in\{+1,-1\}\), the new graph is
\begin{align}
 b^{\rm new}_{i^t h^{u'}}
   &=
   \begin{cases}
     t b_{ih},&t=u',\ i\ne h,\\
     t b_{iP},&t\ne u',
   \end{cases}\notag\\
 b^{\rm new}_{i^t y}&=t b_{iy},\qquad
 b^{\rm new}_{y i^t}=t b_{yi},\qquad
 b^{\rm new}_{y y'}=0\quad(y\ne y'),
 \label{eq:src24}
\end{align}
where \(i,h\in H\) and \(y,y'\in Y\).
Indeed \(P=vH_R/s\) at a left prime and
\(P=-wH_L/s\) at a right prime. Substituting these into the
incoming pivot factor gives the cross-copy rows above.
The left unary acquires
\(\lambda_i(v/s)^{b_{iP}}\), and the right unary acquires
\(\lambda_i(-w/s)^{-b_{iP}}\) after conjugation.
Shared outside rows cancel on other shared outside slots;
their unary factors combine.
For a regular slot \(b_{iP}=\varepsilon_i\), so the old
frequency power cancels and leaves
\(\lambda_i(s)^{-t\varepsilon_i}\), with any sign multiplier
absorbed in \(\kappa_i\). All three invariants follow.
Future active atoms are copied together, so they retain common
incoming columns and total-product dependence.
Thus the off-diagonal part of the extended sum in
\eqref{eq:src22}, before its factor \(e^{d_j}\), is precisely
the next amplitude \(\eta_j\).

\pdfbookmark[2]{Frequency histories and their square weights}{characters.histories}
\subsection*{Frequency histories and their square weights}

Unroll a level-\(l\) term from its root. At a node of step \(j\),
the output slots are \(Y,H_L,H_R\). Each child keeps \(Y\) and
one copy of \(H\), and inserts the pivot integer
\eqref{eq:src23}. Inserted ancestor pivots may occur in later
substitutions. Every newly inserted pivot is required to be
coprime to every frequency below that node. All inherited
supports and ranges are evaluated recursively. We also use the
same histories with one current active pivot replaced by a
fixed external integer, after its outgoing phase has been
removed as in \eqref{eq:src22}.

For a fixed top list assignment, including this external
version, a given root frequency has at most one valid history.
At the top node compare two possibilities
\((v,w,P)\) and \((v',w',P')\). The unit conditions on the
output products give \(s\mid vw'-v'w\), while
\[
 H_R(vw'-v'w)=s(Pw'-P'w).
\]
Consequently,
\[
 |vw'-v'w|
 \le |s|\frac{P|w'|+P'|w|}{H_R}<|s|,
\]
because \(\Delta_j-E_{j-1}\gg\sqrt m\).
The determinant vanishes. The pairs and positive pivots are
proportional. In lowest terms the proportionality numerator
and denominator divide the corresponding frequency and pivot
simultaneously, so the required coprimalities force their ratio
to be 1. This fixes the child frequencies and child lists;
recursion proves uniqueness.

We claim, with \(r=2^l\), that
\begin{equation}\label{eq:src25}
 \E\sum_{\rm histories}|W_l|^2
       \le \exp(Crm+o_k(m)).
\end{equation}
This includes the fixed-external version, averaged over its
other actual priors. Each history has \(r\) bottom weights,
and the explicit level-zero product range gives
\[
 |W_l|^2\le\exp(-r\Delta_0+O_k(1)).
\]
After using this pointwise bound, we may discard the word-bin,
archimedean range, and internal distinctness indicators when
bounding the remaining nonnegative support probability. We retain
the integrality and frequency-unit conditions used below.
There is one independent active word at every bottom leaf.
Fix all frequencies and all actual labels except one chosen
bulk prime in each such word. Let \(R\) be the product of the
absolute node and leaf frequencies, and put
\(Q=R^{k+2}\). Thus \(Q=\exp(O_k(m))\).
The joint residues of the chosen primes modulo \(Q\) are
dominated, for upper bounds by residue conditions, by
independent uniform units at a cost \(\exp(CrL)\).
To see this for one prime, its prior has point weight \(O(1/p)\).
In a unit residue class modulo \(Q\), the sum of \(1/n\) over
each dyadic interval in its range is \(O(1/Q)\), since all these
integers are much larger than \(Q\). There are at most
\(\exp(CL)\) such intervals. The resulting upper bound
\(\exp(CL)/Q\) is no larger than \(\exp(CL)\) times the
uniform unit mass \(1/\varphi(Q)\). Independence proves
the joint assertion.

Under these uniform unit priors, expose first the product of
all chosen variables, then the left-child product at each
split, proceeding downwards. Given the parent product, the
left product is uniform on the unit group and determines the
right product; this is the elementary counting property of
independent uniform group variables. At a node write
\[
 H_L=C_LX_L,\qquad H_R=C_RX_R,
\]
where \(X_L,X_R\) are the products of the chosen variables in
the two child subtrees. The constants can include inserted
ancestor pivots, but their needed residues are known from
previously exposed splits. Every chosen word remains in an
\(H\)-branch through the forward transfers, so a later split
refines an aggregate product already exposed at each ancestor.
The recurrence \eqref{eq:src23} therefore uses the two current
child totals and previously known integers.
After \(h\) reconstructed divisions, those integers remain known
modulo \(R^{k+2-h}\). Indeed, if \(s\mid R\), a numerator known
modulo \(R^a\) determines its integral quotient by \(s\) modulo
\(R^{a-1}\). Products of known residues lose no further precision.
There are at most \(k\) divisions on a branch, leaving enough
precision for every frequency divisibility and unit condition.
A fixed external pivot is known from the outset and consumes
no precision.

On valid support \(C_L,C_R\) are units modulo \(|s|\).
Given their residues and \(X_LX_R\), integrality requires
\[
 vC_RX_R\equiv wC_LX_L\pmod{|s|}.
\]
It is impossible unless \(v,w\) have the same gcd with \(s\).
In the soluble case let \(g=(v,w,s)\). Dividing by \(g\)
reduces the congruence to a unit square-root equation for
\(X_L\) modulo \(|s|/g\). The number of unit roots is at most
\(2^{\omega(|s|/g)+1}\); the elementary divisor and totient
bounds for integers at most \(\exp(O_k(m))\) consequently give
conditional probability at most
\[
 \exp(o_k(m))\frac{g}{|s|}.
\]
The sequential exposure gives the product of these bounds
over the nodes. For fixed child frequencies,
\[
 \sum_{0<|s|\le V_j}\frac{(v,w,s)}{|s|}
 \le 2\sum_{d\mid(v,w)}\ \sum_{h\le V_j/d}\frac1h
 =\exp(o_k(m)).
\]
Sum the internal frequencies from the root down, always using
this bound with their child values fixed. The remaining leaf
choices number at most
\((2V_0)^r=\exp(r\Delta_0+o_k(m))\).
They cancel the preceding square-weight factor.
The remaining cost \(\exp(CrL+o_k(m))\) is at most the
right side of \eqref{eq:src25}. In particular the leading
constant \(C\) can be independent of \(k\).

\pdfbookmark[2]{Cancellation for a one-sided character interaction}{characters.interaction}
\subsection*{Cancellation for a one-sided character interaction}

We need a uniform comparison estimate in two settings:
\begin{enumerate}[label=(\alph*)]
\item two final-list terms related by permutations of bulk slots,
with the same root frequency;
\item two terms on the diagonal of the square in
\eqref{eq:src22}, with a fixed external \(P\), common root
\(v=w\), and a fixed matching identifying the constituent
prime assignments of \(H_L,H_R\), while \(Y\) is shared.
\end{enumerate}
Fix every frequency in both histories and all the matching or
permutation data. In case (b), sample actual variables only on
\(H_L,Y\); the point weight for the matched counterpart will
be factored out below. The phases at actual primes have the
same translating centers, independently of their slots.
Their additive factors therefore cancel in the quotient:
the total prime product and root frequency are unchanged.
On the diagonal the omitted pivot factors are omitted on both
sides, with the same fixed total \(P\).

Suppose the quotient graph has, between a shorter prime \(q\)
and a longer prime \(p\), a factor
\(\xi_q(p)\), where \(\xi_q\) is a nonprincipal character
modulo \(q\), and has no reverse character factor between
this pair. Suppose also that their \(u=\log\log\) ranges have
gap \(\gg L\). We claim that its weighted expectation is
\begin{equation}\label{eq:characters-one-sided}
 O\!\left(\exp(-c_2e^{\alpha L})\right)
\end{equation}
for some \(c_2>0\), uniformly in the fixed data. Bounded unary
restrictions from counterpart priors are permitted, as is the
factor \(e^{T_j+\Delta_j}/H_L\) on the range
\eqref{eq:src20}. We prove the assertion with the support
conditions included.

First consider large coprimalities inside the histories.
At the top, the actual atoms must be pairwise coprime,
internally distinct where required, and coprime to the
external integer when one is present.
At a lower node, pairwise coprime output atoms and
\eqref{eq:src23} make the inserted pivot automatically
coprime to each atom in \(H_L,H_R\). Indeed, a common divisor
with one side would divide its opposite product times \(v\)
or \(w\); the opposite coprimality and frequency-unit
conditions exclude this. For actual primes the latter
conditions follow from their sizes, and for previously
inserted pivots they are retained frequency support conditions.

An inserted ancestor pivot always occupies an active pivot
position. In the earlier forward transfer its designated
active clone was in \(H\), so reversing that transfer places
it in one of \(H_L,H_R\), never in \(Y\). This remains true
at every earlier step and also for a fixed external pivot.
Thus the slots in \(Y\) at any node are actual prime slots.
Apart from the top conditions, the only large coprimalities
not already automatic are between a newly reconstructed
pivot and actual primes in its \(Y\).

These remaining checks can be removed at a cost
\eqref{eq:characters-one-sided}, or the valid support is
empty. To prove this, express every reconstructed integer
by repeated substitution of \eqref{eq:src23}. Its numerator,
after clearing denominators, is a polynomial in the independent
actual prime variables; every denominator is a product of
small frequencies. In the matched setting, an identified
prime is represented by one variable, not by independent
variables for its two occurrences. The independent sampling
variables are precisely those of \(H_L,Y\), with their original
priors. For fixed \(k\), the degrees are bounded by a
polynomial in \(m\), and the logarithms of the absolute
numerator values throughout the ranges are at most
\(\exp(O_k(L))\). The same bound holds after setting a
variable to zero. These facts follow by induction through
the fixed-depth additions and products. A fixed external
integer satisfies \(\log P\le\exp(O_k(L))\), so it respects
the same bounds.

For a required coprimality to a prime variable \(x\), set
\(x=0\) in the numerator polynomial. If the resulting
polynomial is identically zero, the numerator is divisible
by \(x\) for every assignment. Its denominator is a unit
modulo \(x\), so the required coprimality fails identically;
this history pair can be discarded. Otherwise, under
independent sampling of the other actual primes, the
probability that the resulting nonzero polynomial evaluates
to zero is at most its degree times the largest point
mass of any variable. This elementary bound follows
inductively by viewing a nonzero polynomial as a univariate
polynomial, excluding zeros of its leading coefficient,
and using its degree bound on roots. The original independent
priors dominate those with any internal tuple restrictions,
and every actual-prime point mass is at most
\[
 \exp(-c e^{\alpha L}).
\]
On a nonzero evaluation, there are at most
\(\exp(O_k(L))\) possible prime divisors \(x\).
Multiplying their number by the point-mass bound proves
that this coprimality fails with negligible probability.
The same point-mass and divisor argument handles actual
collisions at the top and coprimality to the fixed external
integer. A union bound covers the finitely many checks in
both histories. The factors \(\exp(\operatorname{poly}_k(m))\)
that can multiply these errors remain negligible compared
with \(\exp(c e^{\alpha L})\).

This removal is performed after the additive phases have
canceled and the graph quotient has been written down.
On the enlarged domain its multiplicative factors may use
zero, or any bounded convention, at nonunits; the difference
is supported on the exceptional assignments just bounded.
There is therefore no need to extend an additive inverse
through a failed coprimality.

The remaining integrality and small-frequency gcd conditions
are residue conditions modulo an integer
\[
 Q_*=\exp(O_k(m))
\]
whose prime factors divide the frequencies. One can take
a sufficiently high fixed-for-\(k\) power of their product:
clearing the frequency denominators then determines all
divisibilities and gcds from residues modulo \(Q_*\).
Every actual prime is coprime to \(Q_*\).
The other supports are polynomial inequalities after the
same substitutions: product ranges, pivot ranges, positivity,
and word bins.
Fix all other variables and the short prime. On a dyadic
interval for the long variable, split at the boundaries of
these inequalities and at the critical points of the
polynomial arguments of the smooth factors. Their degrees
and their number are bounded by a polynomial in \(m\) for
fixed \(k\); identically constant polynomials require no
split. On each resulting interval the smooth arguments
are monotone. The level-zero weights
\((X/M)^{1/2}\widehat\psi(vX/M)\), restricted to their stated
\(M\)-ranges, have supremum and variation at most
\(\exp(\operatorname{poly}_k(m))\), by smooth compact support
of \(\widehat\psi\). Products of these weights, and the
optional ratio \(e^{T_j+\Delta_j}/H_L\) on
\eqref{eq:src20}, have the same kind of bound.
Thus, in every allowed residue class, the full remaining
archimedean weight has supremum and total variation at most
\(\exp(\operatorname{poly}_k(m))\).
Individual prime-set membership, cell restrictions, and all
other irregular unary factors can stay in bounded unary
functions.

After fixing the other variables the average consequently
has the form
\[
 \E_p\E_q U(p)V(q)\xi_q(p)W(p,q),
 \qquad |U|,|V|\le1,
\]
with the residue restrictions included in \(W\).
Cauchy--Schwarz removes \(U(p)\). Its resulting nonnegative
square average can be extended from the long-prime prior
to integers using its point bound \(e^{CL}/p\).
Expand the square in \(q,q'\).
The terms \(q=q'\) cost at most the largest short-prime
point mass times \(\exp(\operatorname{poly}_k(m))\).
For \(q\ne q'\), the function
\(\xi_q(n)\overline{\xi_{q'}(n)}\) has mean zero modulo
\(qq'\), since both constituent characters are nonprincipal.
It also has mean zero along a progression of step \(Q_*\),
because \(Q_*\) is coprime to \(qq'\).
On a dyadic interval around \(b\), complete-period
cancellation and partial summation therefore bound the sum,
including all residue classes, by
\[
 \exp(\operatorname{poly}_k(m))\frac{Q_*qq'}{b}.
\]
For example, in one residue class the incomplete character
sum has absolute value at most \(qq'\); multiplying by
the variation of \(W(n,q)\overline{W(n,q')}/n\) gives
the bound without \(Q_*\), and summing the classes gives
the displayed expression.
The gap in \(u\)-ranges means that the smallest long
\(\log b\) exceeds the largest short \(\log q,\log q'\)
by a factor \(\exp(\Omega(L))\). Hence this bound is
exponentially smaller than \(\exp(-c e^{\alpha L})\).
There are only \(\exp(O(L))\) dyadic intervals.
Taking the square root and absorbing the preceding
coprimality errors proves \eqref{eq:characters-one-sided}.
Its strength also permits summing over all fixed-depth
frequency histories and matchings, whose number is at most
\(\exp(\operatorname{poly}_k(m))\).

\pdfbookmark[2]{Anchor codes and the diagonal bound}{characters.anchors}
\subsection*{Anchor codes and the diagonal bound}

An active bulk slot at level \(l\) has a path
\(t=(t_1,\ldots,t_l)\in\{+1,-1\}^l\).
Write \(e_0=1\), \(e_j=\prod_{h\le j}t_h\), and
\(\varepsilon(t)=e_l\). For an anchor clone \(y\) created
by one of the first \(l\) steps, define its code coordinate
at the slot by
\[
 c_y(t)=b_{yi}/\varepsilon(t).
\]
These coordinates depend on the path, not on the bulk
position within a word.
Before step \(j\), a fresh anchor \(a\) has
\[
 b_{ai}=e_{j-1},\qquad b_{aP}=\alpha_j,
\]
where \(\alpha_j\) is the common parity of the active pivot.
Both coefficients initially equal 1, and until its reserved
step the anchor stays outside; the update
\eqref{eq:src24} multiplies an incoming coefficient by the
copy sign of its target.
At its step, the ordered pair from the positive and negative
anchor clones is therefore
\begin{equation}\label{eq:characters-anchor-pair}
 \big(c_{a^+}(t),c_{a^-}(t)\big)=
 \begin{cases}
   (1,-\alpha_j/e_{j-1}),&t_j=+1,\\
   (-\alpha_j/e_{j-1},1),&t_j=-1.
 \end{cases}
\end{equation}
The formula is identical for the small and big anchors.
After their creation the anchors stay outside. Later bulk
copying multiplies both \(b_{yi}\) and \(\varepsilon_i\)
by the same sign, so each coordinate is thereafter fixed.

The pair in \eqref{eq:characters-anchor-pair} is \((1,1)\)
exactly when \(e_{j-1}=-\alpha_j\); either mixed pair
corresponds to \(e_{j-1}=\alpha_j\). It therefore determines
the incoming parity. All anchor pairs together determine
\(e_0,\ldots,e_{l-1}\), hence
\(t_1,\ldots,t_{l-1}\). At most the last sign remains
undetermined. Thus any code occurs on at most two paths,
and code together with final parity determines the entire
path.

Consider the diagonal at step \(j\), put \(l=j-1\), and
write \(r=2^l\). Since every prime factor of \(H_L,H_R\)
exceeds the frequency bounds, the equation
\(vH_R=wH_L\) forces
\[
 H_L=H_R,\qquad v=w.
\]
The products are squarefree on each branch, so equality gives
a unique matching of their constituent slots. Sum by these
matchings. Bulk slots can match only bulk slots because their
size region is separated from every other role in \(H\).
Their original orientations and character rules coincide.

Suppose a matching changes a bulk code. Choose an earlier
anchor coordinate at which the old and new codes differ.
This anchor is in the shared list \(Y\). Write the parities
as \(\varepsilon,\varepsilon'\) and the two coordinates as
\(c,c'\). They are signs, so \(c'=-c\).
The net exponents in the phase quotient from anchor to bulk
and from bulk to anchor are respectively
\[
 \varepsilon c-\varepsilon'c',
 \qquad \varepsilon-\varepsilon'.
\]
If the parities agree, these are \((\pm2,0)\);
use the corresponding small anchor. If the parities differ,
they are \((0,\pm2)\); use the corresponding big anchor.
The small and big anchors have identical code patterns,
so both choices retain the required differing coordinate.
In the first case the short modulus is the small-anchor
prime, and in the second it is the bulk prime.
In either case the surviving interaction is a one-sided
square of a chosen character. It is nonprincipal because
that character has order greater than 2.
All other phase factors are unary once the other variables
and histories are fixed, so
\eqref{eq:characters-one-sided} applies.

The normalization preceding this application is important.
For a fixed assignment on \(H_L\), the point weight of the
matched ordered counterpart is
\[
 \frac{C_H}{H_L}
\]
times its role and distinctness indicators, where
\[
 C_H\le L^{-rm}\exp(Crm+O_k(1)).
\]
There are \(rm\) bulk priors, \(r\) top priors, and only
\(O(n_c)\) cell priors in \(H\); the last contribute
at most \(\exp(Cn_cL)\), which is absorbed uniformly
by the displayed bound. Before summing over the external
integer \(P\), extract
\[
 C_H e^{-T_j-\Delta_j}.
\]
The remaining ratio \(e^{T_j+\Delta_j}/H_L\) is bounded
on \eqref{eq:src20} and was included in the variation
estimate. There are at most \(\exp(T_j+O_k(1))\)
possible \(P\)'s, canceling the large factor \(e^{-T_j}\).
Counterpart role membership is a bounded unary restriction
after the matching; its internal coprimalities are handled
by the cleanup already proved. Thus all code-changing
matchings, all their histories, and the extended sum over
\(P\) have total negligible contribution. This order of
normalization avoids multiplying a small uniform error
by an unbalanced doubly exponential count of \(P\)'s.

For code-preserving matchings there are at most \(2m\)
counterparts for each bulk slot: at most two paths share
its code and each contains \(m\) bulk positions.
There are only boundedly many other slots for fixed \(k\).
Consequently their number is at most
\[
 (2m)^{rm}\exp(O_k(1)).
\]
Fix \(P\) and one such matching. For each common root
frequency \(v\), history uniqueness permits at most one
valid history in either branch. Bound the phases absolutely
and use
\[
 2|W_lW_l'|\le |W_l|^2+|W_l'|^2
\]
on simultaneous support. In each term separately, bound
the point weight of the other branch by
\[
 C_H\exp(-T_j-\Delta_j+O_k(1)),
\]
and retain the square branch's own priors. Equation
\eqref{eq:src25}, including its fixed-external version,
then bounds its frequency sum by
\(\exp(Crm+o_k(m))\). Summing over \(P\) and the matchings,
and using \(m/L=z+o(1)\), bounds this part of the diagonal
in the extended sum, before multiplication by \(e^{d_j}\),
by
\begin{equation}\label{eq:src26}
 \exp\!\left(-\Delta_j+rm(\log z+C)+o_k(m)\right).
\end{equation}

We now specify the order of the constant choices and the
amplitude induction. The construction of
\eqref{eq:src18} fixed \(B_0\) independently of \(B_D,k\).
Choose \(B_1>B_0+2\). Choose \(B_D\) large enough for the
repeat bound and for
\eqref{eq:src26} to be at most
\(\tfrac12 e^{-d_j}e^{-2B_1rm}\), after allowing the
negligible code-changing error.
This is possible uniformly in \(k\): the main negative
term in \eqref{eq:src26} is
\(-rm(B_D+20\log z)\), and its positive term is only
\(rm(\log z+C)\).
Increase \(K_0\) so that
\[
 \sum_{j=1}^k d_j\le Cn_cL+O_k(1)<m
\]
for large \(L\); this follows from
\(n_c/z\ll k e^{-\epsilon k}\).
Whenever \(|\eta_{j-1}|\ge e^{-B_1rm}\), the diagonal
bound and \eqref{eq:src22} give
\[
 |\eta_j|\ge \tfrac12e^{-d_j}|\eta_{j-1}|^2.
\]
To see that the budget is preserved despite the additional
factors, set \(\mathfrak a_j=-2^{-j}m^{-1}\log|\eta_j|\).
The recurrence gives
\[
 \mathfrak a_j\le B_0+
       \sum_{h=1}^j\frac{d_h+\log2}{2^hm}
       \le B_0+1+o_k(1)<B_1.
\]
Starting from \eqref{eq:src19}, induction therefore proves
\begin{equation}\label{eq:characters-amplitude-budget}
 |\eta_l|\ge \exp(-B_1\,2^lm)
 \qquad(0\le l\le k).
\end{equation}
All vanishing errors here are taken only after \(B_D,k\)
are fixed.

\pdfbookmark[2]{Final permutation comparison}{characters.permutations}
\subsection*{Final permutation comparison}

At level \(k\), put \(r=2^k\). For each of the \(m\) bulk
positions, independently permute its \(r\) variables among
paths of the same final parity. There are \(r/2\) paths
of each parity, hence
\[
 J_*=\big[(r/2)!^2\big]^m
\]
such reassignments. All bulk priors are identical and
independent, so these permutations preserve the underlying
measure. The word bins and other supports are part of the
reassigned integrand; they have never conditioned the priors.

Let \(\mu\) be the product of the actual priors and counting
measure on \(0<|s|\le V_k\), of total mass at most \(2V_k\).
For a reassignment \(\pi\), let \(\Phi_\pi\) be the
corresponding sum over internal histories at this top
assignment and root frequency. Prior invariance gives
\[
 \int\Phi_\pi\,d\mu=\eta_k.
\]
History uniqueness and \eqref{eq:src25} give
\[
 \|\Phi_\pi\|_{L^2(\mu)}^2
       \le \exp(Crm+o_k(m)).
\]
For distinct \(\pi,\rho\), some actual bulk variable occupies
different paths of the same final parity. Such paths have
different anchor codes. A small anchor consequently gives
a one-sided nonprincipal square-character interaction in
their phase quotient. The additive phases cancel, since
the same prime product, root frequency, and prime-dependent
centers are used in both assignments. Equation
\eqref{eq:characters-one-sided}, summed over the
fixed-depth frequency data, yields
\[
 |\langle\Phi_\pi,\Phi_\rho\rangle|
   \ll\exp(-c_2e^{\alpha L}),
\]
after decreasing \(c_2>0\) if necessary.
Average the functions \(\Phi_\pi\) and apply
Cauchy--Schwarz against the constant function 1. This gives
\begin{equation}\label{eq:characters-final-comparison}
 |\eta_k|^2\le
 2V_k\left(
   \frac{\exp(Crm+o_k(m))}{J_*}
        +O\!\left(\exp(-c_2e^{\alpha L})\right)\right).
\end{equation}
The factorial estimate and \eqref{eq:src17} imply
\[
 \log J_*\ge rm(k\log2-C),\qquad
 \log V_k=rm(B_D+20\log z)+o_k(m),
 \qquad 20\log z=60\epsilon k.
\]
Because \(\log2-60\epsilon>0\), a sufficiently large
\(K_0\), chosen after \(B_D,B_1\), makes the first term
on the right of \eqref{eq:characters-final-comparison}
smaller than \(\exp(-2B_1rm)\), with a fixed exponential
margin. The second term, including its factor \(V_k\),
is smaller than \(\exp(-Crm)\) for every fixed \(C\)
as \(L\to\infty\) at this fixed \(k\).
This contradicts \eqref{eq:characters-amplitude-budget},
and proves Proposition~\ref{prop:characters}.
\end{proof}

\begin{corollary}\label{cor:mixed-flatness}
Define
\[
 F_p(x)=\frac{\ind_{S_p}(x)-\sigma_p}
                {\sqrt{\sigma_p(1-\sigma_p)}},
 \qquad
 g_p(v)=\sqrt p\,\E_{x\in\F_p}F_p(x)e_p(-vx).
\]
For every fixed \(0<\alpha<\beta\), outside a set of primes
of harmonic mass \(o(L)\) in
\(\alpha L\le\log\log p\le\beta L\), one has uniformly
\begin{equation}\label{eq:src27}
 \sigma_p=\tfrac12+o(1),\qquad
 \max_{\substack{a\in\F_p\\\lambda\ {\rm multiplicative}}}
 \left|\E_{v\in\F_p}g_p(v)\lambda(v)e_p(av)\right|=o(1).
\end{equation}
\end{corollary}

\begin{proof}
Apply Lemma~\ref{lem:collision} at, for example,
\(\log X=\exp((\beta+1)L)\). Its prime cutoff contains
the whole band. The nonnegative imbalance term in
\eqref{eq:src2}, together with
\(\log p\ge e^{\alpha L}\), gives
\[
 \sum_{\alpha L\le\log\log p\le\beta L}\frac1p
 \left(\frac1{\sigma_p}+\frac1{1-\sigma_p}-4\right)
 \ll Le^{-\alpha L}.
\]
Since the summand in parentheses controls
\((\sigma_p-\tfrac12)^2\), this proves balance in harmonic
probability.

If \(\lambda\) is nonprincipal, the Gauss formula gives
\begin{align*}
 \E_v g_p(v)\lambda(v)e_p(av)
 &=\frac{\tau(\lambda)}{p^{3/2}}
           \sum_xF_p(x)\overline{\lambda}(a-x)\\
 &=\frac{\tau(\lambda)}{\sqrt p}\,
     \overline{\lambda}(-1)\,
     \frac{\sigma_p}{\sqrt{\sigma_p(1-\sigma_p)}}
        \E_{x\in S_p}\overline{\lambda}(x-a).
\end{align*}
The first two factors on the last line have modulus 1.
The constant part of \(F_p\) has vanished because the
character is nonprincipal. For quadratic \(\lambda\),
Proposition~\ref{prop:quadratic} and its harmonic-band
consequence give a maximal translated bias tending to
zero in harmonic probability. For order greater than 2,
Proposition~\ref{prop:characters} gives the same conclusion,
already with a maximum over the character.
On the balanced primes the displayed scale factor is
bounded.

For the principal character, \(g_p(0)=0\), and additive
inversion gives exactly
\[
 \E_v g_p(v)\lambda(v)e_p(av)=p^{-1/2}F_p(a).
\]
This tends uniformly to zero on the balanced primes.
Finally choose a positive threshold tending to zero
sufficiently slowly in the balance and maximal-bias
estimates. Markov's inequality makes their combined
exceptional harmonic mass \(o(L)\), proving the asserted
uniform formulation.
\end{proof}
